% 公共导言区：所有章节共用。只在这里新增宏，章节文件里不要 \newcommand。
\usepackage{amsmath,amssymb,amsthm,mathtools}
\usepackage{mathrsfs}
\usepackage{graphicx}
\usepackage{booktabs}
\usepackage{enumitem}
\usepackage{xcolor}
\usepackage[a4paper,margin=2.6cm]{geometry}
\usepackage{tikz-cd}
\usepackage[colorlinks=true,linkcolor=blue!50!black,citecolor=green!40!black,urlcolor=blue!50!black]{hyperref}
\usepackage[capitalize,nameinlink]{cleveref}

% ---------------- 定理环境（按章编号，共用一个计数器）
\theoremstyle{plain}
\newtheorem{theorem}{定理}[chapter]
\newtheorem{proposition}[theorem]{命题}
\newtheorem{lemma}[theorem]{引理}
\newtheorem{corollary}[theorem]{推论}
\theoremstyle{definition}
\newtheorem{definition}[theorem]{定义}
\newtheorem{example}[theorem]{例}
\newtheorem{exercise}{习题}[chapter]
\theoremstyle{remark}
\newtheorem{remark}[theorem]{注}
\newtheorem*{hint}{提示}
\crefname{theorem}{定理}{定理}
\crefname{proposition}{命题}{命题}
\crefname{lemma}{引理}{引理}
\crefname{corollary}{推论}{推论}
\crefname{definition}{定义}{定义}
\crefname{example}{例}{例}
\crefname{exercise}{习题}{习题}
\crefname{remark}{注}{注}
\crefname{chapter}{第}{第}
\crefformat{chapter}{第~#2#1#3~章}
\crefformat{appendix}{附录~#2#1#3}
\crefname{section}{§}{§}
\crefname{equation}{式}{式}
\crefname{figure}{图}{图}
\crefname{table}{表}{表}
\renewcommand{\proofname}{证明}

% 每章开头的「本章要点」框与章末「注记与文献」
\newenvironment{goals}{\par\medskip\noindent\textbf{本章要点}\begin{itemize}[leftmargin=2em,itemsep=1pt]}{\end{itemize}\medskip}
\newcommand{\notes}{\section*{注记与文献}\addcontentsline{toc}{section}{注记与文献}}
\newcommand{\exercises}{\section*{习题}\addcontentsline{toc}{section}{习题}}

% ---------------- 记号（一字一义；见附录 B 记号表）
\newcommand{\ee}{\mathrm{e}}                 % 自然对数底
\newcommand{\ii}{\mathrm{i}}                 % 虚数单位
\newcommand{\dd}{\mathrm{d}}
\newcommand{\R}{\mathbb{R}}
\newcommand{\C}{\mathbb{C}}
\newcommand{\Z}{\mathbb{Z}}
\newcommand{\N}{\mathbb{N}}
\newcommand{\HH}{\mathbb{H}}                 % 上半平面
\newcommand{\D}{\mathbb{D}}                  % 单位圆盘
\newcommand{\id}{\mathrm{id}}
\newcommand{\Rea}{\operatorname{Re}}
\newcommand{\Ima}{\operatorname{Im}}
\newcommand{\Log}{\operatorname{Log}}        % 主值对数
\newcommand{\Phiatt}{\Phi_{\mathrm{att}}}    % 吸引 Fatou 坐标
\newcommand{\Phirep}{\Phi_{\mathrm{rep}}}    % 排斥 Fatou 坐标
\newcommand{\horn}{\mathfrak h}              % horn 映射；逆 horn 映射 \horn^{-1}
\newcommand{\Lam}{\Lambda}                   % 强迫尺度 exp(4π²/log λ₁)
\newcommand{\Ksew}{K^{\mathrm W}}            % 缝合解
\newcommand{\Kkn}{\mathcal K}                % Kneser 解（底数延拓族）
\newcommand{\Pstr}{\mathcal P}               % 抛物层
\newcommand{\Ocal}{\mathcal O}               % 有界全纯函数的 Banach 空间
\newcommand{\Lfun}{\mathfrak L}              % 线性泛函 𝔏_n
\newcommand{\Ifun}{\mathcal I}               % 不变量 𝓘_n
\newcommand{\Dfun}{\mathcal D}               % 抛物层上的对数微分 𝒟_n
\newcommand{\ktr}{\kappa_{\mathrm{tr}}}      % 平移开折的一阶系数
\newcommand{\Lop}{\mathsf L}                 % 同调算子 Lφ = φ∘g − g'φ
\newcommand{\ustar}{u^{*}}                   % 基点（u 坐标）

% ---------------- 页眉：不转大写（避免数学与人名被大写）
\usepackage{fancyhdr}
\setlength{\headheight}{15pt}
\pagestyle{fancy}
\fancyhf{}
\fancyhead[LE,RO]{\thepage}
\fancyhead[RE]{\small\leftmark}
\fancyhead[LO]{\small\rightmark}
\renewcommand{\headrulewidth}{0.4pt}
\fancypagestyle{plain}{\fancyhf{}\fancyfoot[C]{\thepage}\renewcommand{\headrulewidth}{0pt}}
\renewcommand{\chaptermark}[1]{\markboth{\CTEXthechapter\quad #1}{}}
\renewcommand{\sectionmark}[1]{\markright{\thesection\quad #1}}

% 章号用阿拉伯数字（正文引用写作「第~\ref{…}~章」）
\ctexset{chapter={number=\arabic{chapter}}}
